\begin{lemma}[Certified reordered trees in Situation 7.2]
Fix the graph, palette, distinguished edge, ETT, hierarchy
counterexample, minimum context, and a
\noderef{HierarchySituationTwoTailData} \(D\).  Then there is a
\noderef{HierarchySituationTwoReorderedTrees} datum
\(R=(T^-,T^*)\) with the following additional properties.  Put
\[
 n=\mathcal E.\mathit{rounds},\quad
 T_n=\mathcal E.\mathit{series}.\mathit{tree}(n),\quad
 B=\noderef{HierarchyFullTailPrefix}(C,\ell-2).
\]
If \(B=T_n\), then the terminal extension kind is PE.  Each of \(T^-\)
and \(T^*\) has the exact terminal construction route appropriate to
that kind: in the PE case it is obtained from \(T_n\) by a
\noderef{TashkinovAugmentationSequence} under \(\varphi_n\); in the SE
or RE case there are an individual edge \(f_n\), the equality
\(F_n=\{f_n\}\), the exact seed \(T_n+f_n\), and an augmentation
sequence from that seed.  Finally, each reordered tree carries an actual
\noderef{CorrespondingExtendedTashkinovTree} generated by
\(\varphi_n\), and both replacement ETTs have
\noderef{MaximumProperty}.
\end{lemma}
\begin{proof}
Write \(H=C.\mathit{hierarchy}\), let \(q\) be its last divider index,
put \(A=\noderef{HierarchyStarredStage}(H,q)\), and put
\(\ell=\noderef{HierarchyFullTailLength}(C)\).  The tail datum gives
\(\ell\geq2\) and exact indices for \(e_p,y_{p-1},y_p\).  The stage
chain in \noderef{TashkinovHierarchy} gives \(|A|\leq|T|\).  Since
\(\ell=|T|-|A|\),
\[
 |A|+\ell=|T|,\qquad |B|=|T|-2,\qquad |B^+|=|T|-1,           \tag{1}
\]
where \(B=\noderef{HierarchyFullTailPrefix}(C,\ell-2)\) and
\(B^+=\noderef{HierarchyFullTailPrefix}(C,\ell-1)\).
The definition of \noderef{HierarchyFullTailPrefix} is an exact
\noderef{TreeSequencePrefix}; consequently
\(B\preccurlyeq B^+\preccurlyeq T\) retains every displayed vertex and
individual edge identity.  The position equations in \(D\), together
with (1), identify the last two attachment pairs of \(T\) as
\((e_{p-1},y_{p-1})\) and \((e_p,y_p)\).

The attachment field of \noderef{TreeSequence} for \(e_p\) gives
\[
 G.\mathit{ends}(e_p)=
 \{y_p,T.\mathit{vertex}(i)\}\quad\hbox{for some }i<|T|.      \tag{2}
\]
The non-incidence field of \(D\) excludes
\(T.\mathit{vertex}(i)=y_{p-1}\).  Vertex injectivity in
\noderef{TreeSequence} therefore gives \(i<|T|-1\), so the old endpoint
is already in \(B\).  The new endpoint \(y_p\) is absent from \(B\), by
the same injectivity.  Append the individual edge \(e_p\) and \(y_p\)
to \(B\).  Equation (2) verifies attachment; restriction of the old
injective maps, the two vertex exclusions, and edge injectivity verify
the remaining tree-sequence fields.  This gives
\[
                         B+(e_p,y_p)=T^-.                    \tag{3}
\]

The original attachment of \(e_{p-1}\) has its old endpoint below
position \(|T|-1\), hence in \(B\), and its new endpoint is
\(y_{p-1}\).  This vertex and this edge are absent from \(T^-\), by
(3) and the injectivity fields.  Appending them gives
\[
                    T^-+(e_{p-1},y_{p-1})=T^*.              \tag{4}
\]
Restriction of the original display similarly gives
\[
                    B+(e_{p-1},y_{p-1})=B^+,\qquad
                    B^++(e_p,y_p)=T.                        \tag{5}
\]
Equations (3)--(5), with injectivity, give
\[
 V(T^-)=V(T)\setminus\{y_{p-1}\},\quad
 E(T^-)=E(T)\setminus\{e_{p-1}\},
\]
and \(V(T^*)=V(T)\), \(E(T^*)=E(T)\).  Thus these constructions fill
every field of \noderef{HierarchySituationTwoReorderedTrees}; call the
result \(R\).

We extract the certified route to \(B\).  Analyze the terminal
\noderef{TashkinovSeries} certificate using
\noderef{TashkinovSeriesStep}.  In PE it supplies a
\(\varphi_n\)-\noderef{TashkinovClosure} starting at \(T_n\).  In SE or
RE it supplies an exact identity \(f_n\), the equality
\(F_n=\{f_n\}\), the seed \(S=T_n+f_n\), and a closure starting at
\(S\).  The upper-containment field of
\noderef{ExtendedTashkinovTree} places \(T\) in the terminal closure.
Expand the closure and apply
\noderef{TashkinovAugmentationSequenceRestrict} first at \(T\), then at
\(B\).  Since every \noderef{TashkinovAugmentation} raises length by
one, this gives
\[
 T_n\longrightarrow B\quad(\mathrm{PE}),\qquad
 S\longrightarrow B\quad(\mathrm{SE/RE}),                  \tag{6}
\]
except possibly when \(B=T_n\), before the non-PE seed has been added.

We discharge that boundary explicitly.  Suppose \(B=T_n\) and the kind
is SE or RE.  The singleton connecting equality and
\noderef{HierarchyNonparallelFirstEdge} identify the first displayed
edge above \(T_n\) with \(f_n\); (1) and (5) identify it with the
individual edge \(e_{p-1}\).  Hence \(S=B^+\).  The next augmentation
adds \(e_p\), of color \(\theta\), and its missing-color clause gives a
vertex of \(B^+\) missing \(\theta\).  The only new vertex is
\(y_{p-1}\), where \(\theta\) is present by \(D\).  Thus
\[
                       \theta\in\overline\varphi_n(T_n).     \tag{7}
\]

In RE, the current-coloring equality in
\noderef{TashkinovSeriesStep} and its preceding-closed field say that
\(T_n\) is \noderef{ColoredVertexSetClosed} under \(\varphi_n\).  In SE, the preceding tree is
closed and the \noderef{MaximumDefectiveChoice} selected coloring is
\noderef{StableColoring} on \(T_n\).  If a boundary edge had a selected
color missing on \(T_n\), the missing-set clause of stability and the
edge criterion \noderef{sc1} would copy the same colored boundary edge
to the preceding coloring, contradicting preceding closedness.  Hence
\(T_n\) is closed under \(\varphi_n\) in SE too.  But (2) makes \(e_p\)
a boundary edge of \(B=T_n\), while \(D\) gives
\(\varphi_n(e_p)=\theta\) and (7) makes \(\theta\) tree-missing.  This
contradicts closedness.  Therefore
\[
 B=T_n\quad\Longrightarrow\quad\Theta_n=\mathrm{PE}.        \tag{8}
\]
This proves the manuscript's Algorithm 3.1 boundary alternative rather
than invoking TAA without its certified seed.

In the original order, the augmentation \(B\to B^+\) supplies a color
missing on \(B\) for \(e_{p-1}\).  Its witness remains in \(T^-\), so
(4) is a Tashkinov augmentation.  The augmentation \(B^+\to T\)
supplies \(\varphi_n(e_p)=\theta\) and a vertex of \(B^+\) missing
\(\theta\).  Since \(\theta\) is present at the sole new vertex
\(y_{p-1}\), that witness lies in \(B\); hence (3) is also a
Tashkinov augmentation.  Compose these two exact steps with (6), using
(8) in the boundary case.  This gives both asserted PE or exact-seed
SE/RE routes to \(T^-\) and \(T^*\).

The fields of \noderef{HierarchyMinimumCounterexampleContext} say that
\(\varphi_n\) is proper and leaves
exactly \(e_0\) uncolored.  The generating-coloring field of
\noderef{ExtendedTashkinovTree} identifies it with the terminal series
coloring.  The reflexive \noderef{StableColoring} relation on \(T_n\),
with its actual \noderef{TashkinovStoredColors}, therefore satisfies
the hypotheses of \noderef{Td}; thus
\[
                         \noderef{ModColoring}(\mathcal E,\varphi_n).
                                                                    \tag{9}
\]
The exact initial-segment and strict-length clauses above make both
reordered trees eligible for the universal clause of (9), with the
appropriate routes already proved.  It returns an actual
\noderef{CorrespondingExtendedTashkinovTree} for each tree.  Apply
\noderef{MaximumPropertyTransport} separately to these two
certificates and the maximum field of \(C\).  Both replacement ETTs
have maximum property, as required.
\end{proof}
